\section*{提案摘要}
\addcontentsline{toc}{section}{提案摘要}

\subsection*{问题:元火现在的演出流程,卡在哪里}

元火动漫社目前的标准演出布置,大致是这样:

演播室那台机器跑 PPT,导播守在那里。主持人在十几米外的舞台前,手里只有麦克风。这套布置在大部分时间是工作的,但漫展/晚会节奏一紧就有几条隐性的痛:

\begin{itemize}
  \item \textbf{翻页全靠导播自觉}。主持人讲到段落末尾要"下一页",得用麦克风提示导播,或者讲一个明显的过渡词让导播感觉到。导播一走神节奏就脱节。
  \item \textbf{答题环节中转三段}。观众抢答完告诉主持人,主持人对着麦克风复述选项 A/B/C/D,演播室里的导播听见再去鼠标点 PPT 上的按钮。链路里两次人工传话,任何一环卡顿节奏都掉。
  \item \textbf{蓝牙翻页器在元火场地不可用}。主持人和演播室隔十几米加一道布幕,蓝牙/USB 信号到不了。这不是设备问题,是物理问题 —— 任何无线翻页器在这种距离上都得歇菜。
  \item \textbf{主持人看不到时间}。台上手里没有计时器,导播看到时间也没法低声告诉主持人。节目超时是常态。
  \item \textbf{反向通道是空的}。导播能看到全场,但没办法把"还有 2 分钟"递回主持人耳朵里。
\end{itemize}

这套流程的核心问题是单向依赖:演出当中的每一件事,都得从主持人嘴里中转给演播室。主持人和导播之间的沟通鸿沟,是当下大部分节奏问题的根源。

\subsection*{方案:把高频小事的控制权,从演播室搬到主持人手里}

\app\ 不要求改变现有的演播室角色,只重新分配控制权:

\begin{itemize}
  \item \role{导播} —— 还是演播室那台 Windows 机器,装一个 Tauri 桌面应用,看完整画面、管授权、保留紧急锁。
  \item \role{主持人} —— 任何手机,扫一个码进网页,不装 App。
\end{itemize}

主持人手机上能做的事:翻页、触发幻灯片里的互动按钮(答题/抽奖)、申请黑屏/恢复、看自己页面顶部的大字计时。导播随时按一下紧急锁能把所有授权一秒收回。

这套切分有一个明确意图:演出当中的高频小事 —— 翻页、答题按钮、计时 —— 主持人自己掌;少数需要全局视角的事(灯光协调、换 PPT 文件、紧急叫停)还是留在演播室。沟通鸿沟没有被消除,只是从"每一次翻页都靠喊"变成"出意外才需要喊"。

\subsection*{为什么这套不一样}

类似的"网页翻页器"工具不少。\app\ 几处不同:

它的传输有两条路。场地里有靠谱 WiFi 的时候走 LAN 直连,零延迟。但漫展常用场地的实情更接近这样:演播室到舞台十几米,中间隔一道布幕、两堵墙、还有调音台 / 灯光控制器 / 一堆 2.4 GHz 设备的电磁串扰 —— 临时摆的便携路由器不一定撑得住这段物理距离。这时候可以切到架在校园网内部服务器上的中继,我们用 PKU CLAB 验证过:导播机和主持人手机各走各的校园网或外网出口,不再依赖现场局域网信号能不能打通这段距离。主持人手机端代码不变,扫一个不同的二维码而已。

它真的对接 PowerPoint COM,不是模拟翻页键。所以幻灯片里那些点按钮触发 VBA 宏的互动模块,比如答题、抽奖、票选,也能直接接进来 —— 这是普通翻页器和大部分手机翻页 App 干不了的事。

授权链是\,\role{开远控 → 紧急锁 → ARM}\,三道闸串起来,任何一道关上主持人都按不动。这不是为了好看,是为了让导播在演出出岔子时不用先去喊话。

\subsection*{对元火具体能改善什么}

把上面几条痛对应回来:

\textbf{主持人手里有计时器}。主持人页顶部 SLIDE 大字栏就是计时器,超时变红 —— 主持人不需要回头看时钟。"还有 X 分钟" 这种导播没法低声递过去的信息,现在主持人自己看得见。这一条直接消掉了"节目超时"的最大单点原因。

\textbf{答题环节砍掉口头中转}。主持人页中间会自动列出当前幻灯片上所有绑了 VBA 宏的按钮(选项 A / 选项 B / 公布答案)。观众抢答后,主持人自己点哪个就触发哪个,不再"主持人复述 → 导播按"这两段传话。原本要三个人配合的环节,变成主持人一个动作。

\textbf{翻页不依赖导播自觉}。主持人按一下大字"下一页"按钮,PPT 立刻翻 —— 导播是否在看屏幕、是否听到口头提示,都不再影响节奏。导播这边变成只看不动,可以释放一名工作人员的注意力去管别的事。

\textbf{物理信号问题被绕开}。WiFi 在大部分漫展场地是有的,没有的话便携路由器或手机热点也能撑住;主持人和演播室距离 100 米也没问题。这是把"信号能不能到"的物理问题,换成了"网络能不能联"的工程问题 —— 后者好解决得多。

\textbf{出岔子时有结构性的兜底}。漫展节目交接、设备脱手、临时换主持人这些瞬间,导播按一下紧急锁,主持人那边所有按钮一起失效。有这个按钮和没有,事故面差一个量级。

\subsection*{建议}

如果元火愿意拿一场演出试一下,这边可以这样配合:

\begin{enumerate}
  \item 提供安装包,演出当天到现场陪测一次,有问题当场处置。
  \item 之后按你们用下来的反馈做调整 —— 比如做元火专用答题模板、cosplay 评分模板这类常用件。
  \item 校园网 Relay 通道已经就绪(\,PKU CLAB \,),不需要社团这边做任何服务器配置。
  \item 事后在白皮书"致谢"页公开标注合作,就一行字,不是营销文案那种。
\end{enumerate}
